% Candidate section. Read with the active manuscript's existing preamble.
% Mathematical purpose: a complete, non-Lipschitz critical-exponent example.
\section{法向收缩与沿零集的 Hölder 漂移}
\label{sec:triangular-holder-example}

下面考虑一个具有线性法向阻尼与平方根切向漂移的二维三角系统。
一个坐标向零集作线性收缩，另一个坐标沿零集漂移。
这个例子说明，近端映射虽不局部 Lipschitz，切向位移的可和性仍能保证点收敛。

在 $X=\mathbb R^2$ 上定义
\[
 F(u,v)=(u,-\sqrt{|u|}),\qquad
 S=F^{-1}(0)=\{0\}\times\mathbb R,
 \qquad \lambda=1.
\]
近端方程 $(a,b)=(u,v)+F(u,v)$ 可直接求解，因而完整近端映射为
\begin{equation}
\label{eq:triangular-complete-prox}
 J_F(a,b)=\left(\frac a2,\ b+\sqrt{\frac{|a|}{2}}\right),
 \qquad
 C(a,b):=(2J_F-I)(a,b)=\left(0,b+\sqrt{2|a|}\right).
\end{equation}
它在全空间处处存在且唯一。

\begin{proposition}
\label{prop:triangular-critical-rleb}
取 $U=X$、$\mathcal J=J_F$，令 $U_R=\{(a,b):|a|\leq R\}$。
定理~\ref{thm:selection-RLEB} 的全部前件成立，可取
\[
 R=\frac1{16},\qquad \gamma=\frac12,
 \qquad L=\sqrt2,\qquad
 \bar t=1,\qquad \psi(t)=t^2,\qquad \kappa=\frac34.
\]
在零线任一点的邻域，Hölder 指数 $1/2$ 与真实残差误差界幂次 $2$
均为最大可取值。
\end{proposition}

\begin{proof}
完整近端覆盖由 \eqref{eq:triangular-complete-prox} 给出。
对输入 $(a,b)$，最近零点为 $p=(0,b)$，且
\[
 \|C(a,b)-p\|^2=2|a|=L^2d((a,b),S)^{2\gamma}.
\]
因此主定理的锚定条件成立，包括零输入距离的情形。

由于 $F$ 单值，其真实残差满足
\[
 r_F(u,v)^2=u^2+|u|,\qquad
 d((u,v),S)=|u|\leq r_F(u,v)^2.
\]
这一误差界对全部输出成立。评价窗和直接兼容条件分别为
\[
 \begin{aligned}
 \frac{R+L\sqrt R}{2}&=\frac{1+4\sqrt2}{32}<1,\\
 \frac{\psi((r+L\sqrt r)/2)}r
 &=\frac{(\sqrt r+\sqrt2)^2}{4}
 \leq\frac{(1/4+\sqrt2)^2}{4}<\frac34
 \quad(0<r\leq R).
 \end{aligned}
\]
因此 (B3)--(B4) 成立。
$U=X$ 使主定理的留域预算自动成立。

此外，对输入对距 $\delta\leq R$，有全对估计
\[
 \|C(a,b)-C(\widetilde a,\widetilde b)\|
 \leq\delta+\sqrt{2\delta}
 \leq(\sqrt2+\sqrt R)\delta^{1/2}.
\]
这里的全对常数与上述锚定常数不同。
而 $\|C(t,0)-C(0,0)\|=\sqrt{2t}$（$t>0$），
故任意大于 $1/2$ 的局部 Hölder 指数均不成立。
若真实残差误差界的幂次 $q>2$，则
\[
 \frac{d((t,0),S)}{r_F(t,0)^q}
 =\frac{t}{(t+t^2)^{q/2}}\longrightarrow+\infty.
\]
因此幂次 $2$ 也不能增大。
\end{proof}

取本例的锐误差界幂次 $q=2$，则 $\gamma q=1$。
非 Lipschitz 行为发生在沿零集的位移上，并未破坏法向收缩。
这一点还可由轨道的闭式看得更清楚。对任意初值 $(a_0,b_0)$，
\[
 a_k=2^{-k}a_0,\qquad
 b_k=b_0+\frac{\sqrt{|a_0|}}{\sqrt2-1}(1-2^{-k/2}).
\]
故 $d((a_k,b_k),S)=2^{-k}|a_0|$，且
\[
 (a_k,b_k)\longrightarrow
 \left(0,b_0+\frac{\sqrt{|a_0|}}{\sqrt2-1}\right).
\]
步长可以逐项求和；例如轨道总长度不超过
$|a_0|+\sqrt{|a_0|}/(\sqrt2-1)$。
当 $a_0\ne0$ 时，点误差的渐近因子为 $2^{-1/2}$，
而距离零集的因子为 $1/2$。
因此，控制距离并不能省去控制切向漂移的步骤。
